\section{Conventions and reduction to bounded integral distances}
\label{pre:section}

For a nonempty set $H\subseteq F$, write
\[
 d(p,H)=\min_{x\in H}d(p,x),\qquad
 \operatorname{cost}(H)=\sum_{p\in D}d(p,H),\qquad
 \operatorname{opt}_r=\min_{\substack{H\subseteq F\\1\le|H|\le r}}
                          \operatorname{cost}(H).
\]
We use $n=|D|$ and $N=n+|F|$. We use $\log(N+2)$ for logarithmic
factors in the instance size. Constants may depend on fixed accuracy
parameters, but not on the input. The expression $\exp(O(m))$ always
has an input-independent constant in the exponent.

We assume $F\ne\varnothing$ and $k\ge1$. Empty client sets have cost zero.
When $k\ge |F|$, opening every facility is optimal. When $k\ge n$, choose
a nearest available facility for each client; their union has at most
$n$ elements and attains the sum of the individual lower bounds.
We may consequently restrict the nontrivial analysis to
\begin{equation}\label{pre:nontrivial}
 1\le k<n,\qquad k\le |F|.
\end{equation}
If a set has fewer than a desired number of facilities, it can be padded
from $F$ without increasing cost. Whenever exact cardinalities are used
for comparison, this padding is understood and will be specified at
the relevant construction.

Fix nearest assignments whenever a solution is considered. For a
comparison solution $O$ of size $h$, its cluster at $i\in O$ is denoted
$C_i$, and
\[
 n_i=|C_i|,\qquad d_p^*=d(p,O),\qquad
 P_i=\sum_{p\in C_i}d_p^*,\qquad P=\operatorname{cost}(O).
\]
Clusters may be empty. A cluster mean $\bar d_i=P_i/n_i$ is used only
when $n_i>0$. A local solution is denoted $S$, with $s_p=d(p,S)$.
Assignment ties are broken by fixed priorities. When the reference
set is enlarged only by adding points, its old points retain their
relative priority. Thus the clients assigned to any fixed subset of
old points can only disappear from that subset.

The body of the proof first treats disjoint indexed client and facility
sets with positive integer distances between distinct points, bounded
by a fixed polynomial in $N$. In particular every nonempty comparison
solution has integral cost at least $n$. The polynomial bound may have
a fixed coefficient depending on an accuracy parameter. It is fixed
before the surplus algorithm is invoked.

\begin{lemma}[Metric rounding and transfer]\label{pre:rounding}
Fix rational $0<\varepsilon_0\le1$. Suppose a polynomial-time algorithm for the bounded
integral instances just described always returns at most $k$ facilities
of cost at most $5\operatorname{opt}_k$, and, with probability at least
$1-p$, returns cost at most $A\operatorname{opt}_k$, where $A\le5$.
Then a randomized polynomial-time algorithm for arbitrary finite
rational metric instances always returns at most $k$ facilities of
cost at most $5(1+\varepsilon_0)\operatorname{opt}_k$ and, with
probability at least $1-p$, returns cost at most
$A(1+\varepsilon_0)\operatorname{opt}_k$.
The bounded integral distances needed are at most
$\lceil20n^2/\varepsilon_0\rceil+1$.
\end{lemma}

\begin{proof}
Detect zero optimum first. It occurs exactly when every client location
contains an available facility and the number of distinct client
locations is at most $k$. More generally, if repeated indexed points
are initially represented by zero distances, use their zero-distance
equivalence classes for this test. Opening one facility in each relevant
class gives cost zero.

Suppose the optimum cost $P^*$ is positive. Guess a positive client--facility
distance $M$ from the input; there are at most $n|F|$ possibilities.
For the correct guess, $M$ is the largest assigned distance in some
fixed optimum, so
\begin{equation}\label{pre:M-range}
 M\le P^*\le nM.
\end{equation}
Set
\[
 u=\frac{\varepsilon_0M}{n},\qquad
 C=\lceil20n^2/\varepsilon_0\rceil+1.
\]
Separate client and facility copies when a point serves both roles.
For distinct indexed points $x,y$ set
\[
 d_M(x,y)=\min\{C,\max\{1,\lceil d(x,y)/u\rceil\}\},
 \qquad d_M(x,x)=0.
\]
This is a metric. Indeed, ceiling is monotone and subadditive on
nonnegative arguments; enforcing the lower bound one on distinct
points and applying a common cap preserve the triangle inequality.
For a triangle with repeated indices the assertion is immediate.
For three distinct indices it follows by applying these two monotone,
subadditive operations to the original triangle inequality. These
observations also cover distinct copies whose original distance is zero.

The original optimum, evaluated in $d_M$, has cost at most $P^*/u+n$.
Thus for the correct guess the new optimum $Q^*$ satisfies
\begin{equation}\label{pre:rounded-optimum}
 Q^*\le P^*/u+n\le n^2/\varepsilon_0+n.
\end{equation}
The assumed algorithm's output has rounded cost at most $5Q^*<C$.
Choose nearest rounded assignments for that output. None of their
edges has rounded length $C$, because the whole sum is smaller than
$C$. On every selected edge, therefore,
$d(x,y)\le u\,d_M(x,y)$. The original nearest-assignment cost is no
larger than that of these assignments. The output consequently has
original cost at most
\[
 5uQ^*\le5(P^*+un)\le5(1+\varepsilon_0)P^*,
\]
using Equation~\eqref{pre:M-range}. On the successful event the same
argument replaces $5$ by $A$.

Run the algorithm for every positive guess and return the feasible
output with the least original cost. It is at least as good as the
output of the correct guess, so no union bound over guesses is needed.
All sets returned respect the budget. Sorting distances, taking ceilings,
and constructing the rounded matrices require polynomial-bit rational
arithmetic. The number of guesses is polynomial, and the rounded
distance bound is a fixed polynomial for fixed $\varepsilon_0$.
\end{proof}

The factor-five requirement in Lemma~\ref{pre:rounding} is supplied by
always retaining a local-search fallback; see Lemma~\ref{loc:swap}.
The algorithm eventually returns the cheapest of its feasible candidates.
This ensures that failure of a sampling event never causes unbounded cost.
